\documentclass[10pt,a4paper]{amsart}
\usepackage[margin=19mm]{geometry}
\usepackage{amsmath,amssymb,mathtools}
\usepackage[scr=rsfs,cal=euler]{mathalfa}
\usepackage{xfrac}
\usepackage[unicode,hidelinks,bookmarks=false]{hyperref}
\newcommand{\ie}{i.e.\ }
\newcommand{\cf}{cf.\ }
\newcommand{\resp}{respectively\ }
\newcommand{\Subsection}{\subsection}
\providecommand{\nobreakdash}{\nobreak\hbox{-}}
\setlength{\emergencystretch}{2em}
\multlinegap0pt
\usepackage{amssymb}
\usepackage{mathtools}
\usepackage{enumerate}
\usepackage{xcolor}
\newtheorem{lemma}{Lemma}
\title{Additive Bases from Primitive Dyck Words:\\ Regular Underapproximations, Motzkin Coding, and Digit Lifting}
\author{Takayuki Kuriyama}
\date{Source: arXiv:2607.23521v1 (CC BY 4.0)}
\begin{document}
\maketitle

\noindent\emph{Source and licence.} Additive Bases from Primitive Dyck Words:\\ Regular Underapproximations, Motzkin Coding, and Digit Lifting. Version arXiv:2607.23521v1 (CC BY 4.0), distributed by arXiv under the Creative Commons Attribution 4.0 International licence, http://creativecommons.org/licenses/by/4.0/, as recorded in the arXiv licence metadata for this exact version and shown on the abstract page https://arxiv.org/abs/2607.23521v1. Full licence notice: \texttt{licenses/arxiv-2607.23521-CC-BY-4.0.txt}.\par\smallskip

\noindent\textit{Editorial scope.} Counted unit: the article's lemma lem:extremal (source0.tex lines 404--419). The article's complete original proof of it is delivered with it (source0.tex lines 421--506, one \texttt{proof} environment of 86 lines). No mathematics is written, repaired or re-derived: the statement and the proof are the article's own lines, in the article's own notation, and no retained line is substituted or re-flowed.  No further source-local context is needed: every cross-reference of every retained line resolves inside the two delivered spans.  Nothing was omitted from any retained span, so the package carries no omission record.  No retained line contains a citation, so the wrapper carries no bibliography and no bibliography entry.  No retained line contains a figure environment, an image-inclusion command or drawing code, so no image is delivered and the package declares no asset dependency.  Editorial formatting only, in addition: an AMS \texttt{amsart} preamble that restates the article's own declarations and macros used by the retained lines, namely the environments \texttt{lemma} on separate counters, together with the standard packages those lines need.  The retained environments print in this document as Lemma 1 in order of appearance, whereas the article prints the same items with the numbering of its own full document.  The complete native source of the article as distributed in the arXiv e-print for this exact version is bundled unchanged as \texttt{sources/206034/source0.tex}.
\begin{lemma}[Extremal two-coloured Motzkin values]\label{lem:extremal}
For every $\ell\geq0$ and every two-coloured Motzkin word $w$ of length $\ell$,
\begin{equation}\label{eq:motzkin-extrema}
  \frac{4^{\ell}-1}{3}\leq[w]_4\leq4^{\ell}-2^{\ell}.
\end{equation}
Both bounds are attained.  The lower bound is attained by $1^{\ell}$, and the
upper bound is attained by
\[
  w_{\max,\ell}=
  \begin{cases}
    3^a0^a,&\ell=2a,\\
    3^a2\,0^a,&\ell=2a+1.
  \end{cases}
\]
Here $1^0$ is understood to be the empty word $\lambda$.
\end{lemma}

\begin{proof}
We use induction on $\ell$.  The assertion is immediate for $\ell=0$.
Write
\[
 w=c_0c_1\cdots c_{\ell-1},
 \qquad
 c_j\in\{0,1,2,3\}.
\]

\emph{Lower bound.}
Since $\sigma(0)=-1$, the first digit cannot be $0$.  If $c_0\in\{1,2\}$, then $\sigma(c_0)=0$, so deleting the first
digit does not change any subsequent prefix weight or the total weight.
Hence the suffix $w'=c_1\cdots c_{\ell-1}$ is again a two-coloured
Motzkin word.  By the induction hypothesis for words of length $\ell-1$,
\[
  [w]_4\geq4^{\ell-1}+[w']_4
  \geq4^{\ell-1}+\frac{4^{\ell-1}-1}{3}
  =\frac{4^{\ell}-1}{3}.
\]
If $c_0=3$, then
\[
  [w]_4
  \geq[3\underbrace{00\cdots0}_{\ell-1\text{ digits}}]_4
  =3\cdot4^{\ell-1}
  \geq\frac{4^{\ell}-1}{3}.
\]
Equality is attained by $w=1^{\ell}$.

\emph{Upper bound.}
If $c_0\in\{1,2\}$, then, as above, $w'$ is again a two-coloured
Motzkin word.  By the induction hypothesis for words of length $\ell-1$,
\begin{align*}
  [w]_4
  &=[c_0c_1\cdots c_{\ell-1}]_4\\
  &\leq[2c_1\cdots c_{\ell-1}]_4\\
  &\leq2\cdot4^{\ell-1}+\bigl(4^{\ell-1}-2^{\ell-1}\bigr)\\
  &=3\cdot4^{\ell-1}-2^{\ell-1}\\
  &\leq4^{\ell}-2^{\ell}.
\end{align*}

Now suppose $c_0=3$, and let $r$ be the length of the shortest nonempty
prefix of $w$ having total weight $0$.  This prefix has the form $3u0$,
where
\[
 u=c_1\cdots c_{r-2},
\]
and let
\[
 v=c_r\cdots c_{\ell-1}
\]
be the remaining suffix.  Then
\[
 2\leq r\leq\ell,
 \qquad
 |u|=r-2,
 \qquad
 |v|=\ell-r,
\]
and $u,v\in\mathsf{Motz}$.  Thus $w=3u0v$.  By place value and the induction hypothesis,
\begin{align*}
  [w]_4
  &=3\cdot4^{\ell-1}+[u]_4\,4^{\ell-r+1}+[v]_4\\
  &\leq3\cdot4^{\ell-1}
    +\bigl(4^{r-2}-2^{r-2}\bigr)4^{\ell-r+1}
    +4^{\ell-r}-2^{\ell-r}\\
  &=4^{\ell}-2^{2\ell-r}+2^{2\ell-2r}-2^{\ell-r}\\
  &=4^{\ell}-\bigl(2^{2\ell-r}-2^{2\ell-2r}\bigr)-2^{\ell-r}.
\end{align*}
It remains to verify
\[
  2^{2\ell-r}-2^{2\ell-2r}
  \geq2^{\ell}-2^{\ell-r}.
\]
After factoring, this is
\[
  2^{2\ell-2r}(2^r-1)\geq2^{\ell-r}(2^r-1),
\]
which follows from $r\leq\ell$.  This proves the upper bound.

Finally, the displayed words $w_{\max,\ell}$ are two-coloured Motzkin words.
A geometric-series calculation gives
\[
  [w_{\max,\ell}]_4=4^{\ell}-2^{\ell},
\]
so the upper bound is attained.
\end{proof}

\end{document}
